\chapter{Compulsions}

\section*{Definition}
Compulsions are repetitive behaviors (e.g., hand washing, checking, ordering) or mental acts (e.g., counting, praying, repeating words) that an individual feels driven to perform in response to an obsession or according to rigid rules. The behaviors are aimed at preventing or reducing distress or a dreaded event but are not realistically connected to the event.

\section*{Pathophysiology}
Compulsions arise from hyperactivation of the orbitofrontal--anterior cingulate--striatal--thalamic circuitry, with impaired inhibitory control from the prefrontal cortex. Striatal dysfunction (particularly the caudate nucleus) disrupts habitual learning and behavioral inhibition, causing stereotyped, repetitive actions.

\section*{Causes}
\begin{itemize}
  \item \textbf{Psychiatric:} Obsessive-compulsive disorder (OCD), body dysmorphic disorder, hoarding disorder, trichotillomania, excoriation disorder
  \item \textbf{Neurological:} Basal ganglia disorders (Tourette syndrome, Sydenham chorea), Parkinson disease, PANDAS/PANS
  \item \textbf{Genetic:} Heritability ~40--50\%; associated with variants in SLC1A1, GRIN2B, and SAPAP3 genes
  \item \textbf{Psychological:} Learned associations (obsession-relief cycle), inflated responsibility, perfectionism
\end{itemize}

\section*{When to See a Doctor}
See your doctor if:
\begin{itemize}
  \item Compulsions consume more than 1 hour per day or cause marked distress and functional impairment
  \item The individual cannot resist performing compulsions despite recognizing irrationality
  \item Skin picking or hair pulling causes visible tissue damage or infection
\end{itemize}

\section*{Related Symptoms}
\begin{itemize}
  \item Obsessions, anxiety, avoidance behaviors, tics, hoarding, body-focused repetitive behaviors, depression
\end{itemize}
\textbf{Important:} Compulsions are performed to neutralize obsessions but paradoxically reinforce them through negative reinforcement.

\section*{Self-Care / Home Management}
\begin{itemize}
  \item Practice gradually delaying or modifying compulsive responses (e.g., wash 1 minute less each day)
  \item Keep a log of triggers, the compulsion performed, and the distress level (0--10)
  \item Use the "15-minute rule": wait 15 minutes before performing a compulsion and observe if the urge decreases
  \item Engage in competing responses (e.g., fist clenching instead of hair pulling)
\end{itemize}

\section*{How Your Doctor Will Evaluate You}
\begin{itemize}
  \item Clinical interview and DSM-5 criteria; differentiate from tics, stereotypies, and addictive behaviors
  \item Severity: Yale-Brown Obsessive-Compulsive Scale (Y-BOCS) or Children's Y-BOCS
  \item Assess for comorbid tic disorders, attention-deficit/hyperactivity disorder, and depression
  \item In children: consider PANDAS/PANS if acute onset following streptococcal infection
\end{itemize}

\section*{Treatment Options}
\begin{itemize}
  \item First-line: Cognitive-behavioral therapy with exposure and response prevention
  \item First-line pharmacotherapy: SSRIs (fluoxetine, fluvoxamine, sertraline) at higher doses than for depression
  \item Augmentation with antipsychotics (aripiprazole, risperidone) for treatment-resistant cases
  \item Severe, refractory OCD may benefit from deep brain stimulation or ablative neurosurgery
\end{itemize}

\clearpage
